\documentclass[a4paper,12pt]{article}

\usepackage[T2A]{fontenc}
\usepackage[utf8]{inputenc}
\usepackage[russian]{babel}
\usepackage{amsmath,amssymb,mathtools}
\usepackage{PTSans}
\renewcommand{\familydefault}{\sfdefault}
\usepackage{icomma}
\usepackage[a4paper,margin=2.05cm,headheight=15pt]{geometry}
\usepackage[nopatch=footnote]{microtype}
\usepackage{xcolor}
\usepackage{tikz}
\usetikzlibrary{calc,arrows.meta,positioning,shapes.geometric}
\usepackage{tabularx,booktabs}
\usepackage{enumitem}
\usepackage{hyperref}
\usepackage{parskip}
\usepackage{fancyhdr}
\usepackage{setspace}
\usepackage[most]{tcolorbox}
\usepackage{titlesec}

\definecolor{accent}{HTML}{008080}
\definecolor{darkaccent}{HTML}{004D4D}
\definecolor{textgray}{HTML}{333333}
\definecolor{softaccent}{HTML}{EAF5F5}
\definecolor{softgray}{HTML}{F5F6F6}

\hypersetup{
  colorlinks=true,
  linkcolor=darkaccent,
  urlcolor=darkaccent
}

\color{textgray}
\onehalfspacing
\setlength{\parindent}{0pt}

\setlist[itemize]{
  leftmargin=1.5em,
  itemsep=0.25em,
  topsep=0.35em
}

\setlist[enumerate]{
  leftmargin=1.65em,
  itemsep=0.35em,
  topsep=0.35em
}

\titleformat{\section}
  {\Large\bfseries\color{darkaccent}}
  {\thesection.}{0.55em}{}

\titleformat{\subsection}
  {\large\bfseries\color{accent}}
  {\thesubsection.}{0.5em}{}

\titlespacing*{\section}{0pt}{1.2em}{0.55em}
\titlespacing*{\subsection}{0pt}{0.9em}{0.4em}

\pagestyle{fancy}
\fancyhf{}
\fancyhead[L]{\small Тригонометрические уравнения}
\fancyhead[R]{\small Тимофей Васильченко}
\fancyfoot[C]{\small\thepage}
\renewcommand{\headrulewidth}{0.4pt}
\renewcommand{\footrulewidth}{0pt}

\newtcolorbox{keybox}[1][]{
  enhanced,
  breakable,
  colback=softaccent,
  colframe=accent!55,
  boxrule=0.5pt,
  arc=1.5mm,
  left=3mm,
  right=3mm,
  top=2.5mm,
  bottom=2.5mm,
  before skip=0.7em,
  after skip=0.7em,
  #1
}

\newtcolorbox{noteBox}[1][]{
  enhanced,
  breakable,
  colback=softgray,
  colframe=black!15,
  boxrule=0.4pt,
  arc=1.5mm,
  left=3mm,
  right=3mm,
  top=2.5mm,
  bottom=2.5mm,
  before skip=0.7em,
  after skip=0.7em,
  #1
}

\tikzset{
  flowStart/.style={
    ellipse,
    draw=darkaccent,
    fill=softaccent,
    minimum width=4.6cm,
    minimum height=1.05cm,
    align=center,
    font=\small
  },
  flowProcess/.style={
    rectangle,
    rounded corners=2mm,
    draw=accent!75!black,
    fill=white,
    minimum width=8.4cm,
    text width=7.7cm,
    minimum height=1.05cm,
    align=center,
    font=\small,
    inner sep=3mm
  },
  flowDecision/.style={
    diamond,
    aspect=2.5,
    draw=darkaccent,
    fill=softaccent,
    text width=4.9cm,
    align=center,
    font=\small,
    inner sep=1.2mm
  },
  flowArrow/.style={
    -{Stealth[length=2.8mm,width=1.8mm]},
    line width=0.8pt,
    draw=darkaccent
  }
}

\newcommand{\Z}{\mathbb{Z}}

\begin{document}

% ============================================================
% Титульный лист
% ============================================================

\hypersetup{pageanchor=false}

\begin{titlepage}
\thispagestyle{empty}

\begin{tikzpicture}[remember picture,overlay]
  \draw[
    accent!45,
    line width=1.25pt
  ]
    ($(current page.south west)+(0,1.55cm)$)
    .. controls
    ($(current page.south west)+(5.2cm,2.45cm)$)
    and
    ($(current page.south east)+(-6.0cm,0.85cm)$)
    ..
    ($(current page.south east)+(0,1.75cm)$);

  \draw[
    darkaccent!25,
    line width=0.65pt
  ]
    ($(current page.south west)+(0,1.15cm)$)
    .. controls
    ($(current page.south west)+(6.8cm,1.75cm)$)
    and
    ($(current page.south east)+(-4.8cm,1.05cm)$)
    ..
    ($(current page.south east)+(0,1.25cm)$);
\end{tikzpicture}

\vspace*{4.2cm}

\begin{center}

{\large ЕГЭ по профильной математике\par}

\vspace{1cm}

{\fontsize{28}{34}\selectfont
\bfseries
\color{darkaccent}
Чек-лист\par}

\vspace{0.55cm}

{\LARGE\bfseries
Тригонометрические уравнения\par}

\vspace{0.25cm}

{\Large
с радикалами и отбором корней\par}

\vspace{1.8cm}

{\large Тимофей Васильченко\par}

\vspace{0.45cm}

{\large \today\par}

\vfill

{\normalsize
Лёвин Артём Александрович эксклюзивно для Тимофея Васильченко\par}

\vspace{1.8cm}

\end{center}

\end{titlepage}

\hypersetup{pageanchor=true}

% ============================================================
% Основной чек-лист
% ============================================================

\section{Главный равносильный переход: уравнение с корнем}

Для уравнения вида
\[
\sqrt{F(x)}=G(x)
\]
правильный равносильный переход имеет вид
\[
\boxed{
\sqrt{F(x)}=G(x)
\iff
\begin{cases}
F(x)=G^2(x),\\
G(x)\ge 0.
\end{cases}}
\]

\begin{keybox}
\textbf{Почему отдельное ограничение на подкоренное выражение здесь не требуется?}

После условия
\[
F(x)=G^2(x)
\]
автоматически получаем
\[
F(x)\ge0.
\]
Поэтому при таком равносильном переходе критично сохранить условие
\[
G(x)\ge0.
\]
\end{keybox}

Для уравнения с двумя корнями:
\[
\sqrt{F(x)}=\sqrt{G(x)}
\iff
\begin{cases}
F(x)=G(x),\\
F(x)\ge0.
\end{cases}
\]

После равенства \(F(x)=G(x)\) вместо условия \(F(x)\ge0\)
можно использовать условие \(G(x)\ge0\).

\subsection{Ключевая проверка перед возведением в квадрат}

\begin{itemize}
  \item Определите, какая часть уравнения обязана быть неотрицательной.
  \item Запишите ограничение до дальнейших преобразований.
  \item После получения тригонометрических корней вернитесь к ограничению и отфильтруйте неподходящие серии.
\end{itemize}

\section{Два ориентира при преобразовании тригонометрии}

При сложном уравнении полезно последовательно задать два вопроса.

\begin{enumerate}

  \item
  \textbf{Одинаковы ли аргументы?}

  Если встречаются
  \[
  x,\qquad
  \frac{\pi}{2}+x,\qquad
  \pi-x
  \]
  и другие различные аргументы, сначала примените формулы приведения.

  \item
  \textbf{Можно ли перейти к одной тригонометрической функции?}

  Если присутствуют квадраты
  \[
  \sin^2x,\qquad \cos^2x,
  \]
  используйте основное тригонометрическое тождество.

\end{enumerate}

\begin{keybox}
\textbf{Стратегическая цель:}
привести выражение к алгебраическому уравнению относительно одной величины, например
\[
t=\cos x
\qquad\text{или}\qquad
t=\sin x.
\]
\end{keybox}

Полезные формулы приведения:
\[
\cos\left(\frac{\pi}{2}+x\right)=-\sin x,
\qquad
\cos\left(\frac{\pi}{2}-x\right)=\sin x,
\]
\[
\sin\left(\frac{\pi}{2}+x\right)=\cos x,
\qquad
\sin\left(\frac{\pi}{2}-x\right)=\cos x.
\]

Основное тригонометрическое тождество:
\[
\sin^2x+\cos^2x=1.
\]

Следствия:
\[
\sin^2x=1-\cos^2x,
\qquad
\cos^2x=1-\sin^2x.
\]

\begin{noteBox}
\textbf{Правило скобок.}
Подставляемое выражение берите в скобки.

Например:
\[
-4\cos^2x
=
-4\bigl(1-\sin^2x\bigr).
\]

Так сохраняются все знаки при раскрытии скобок.
\end{noteBox}

\section{Разобранный образец}

Рассмотрим уравнение
\[
\sqrt{2\cos^3x-2\cos x}
=
\cos\left(\frac{\pi}{2}+x\right).
\]

По формуле приведения
\[
\cos\left(\frac{\pi}{2}+x\right)
=
-\sin x.
\]

Получаем равносильную систему
\[
\begin{cases}
2\cos^3x-2\cos x=\sin^2x,\\
-\sin x\ge0.
\end{cases}
\]

Ограничение:
\[
\sin x\le0.
\]

Используем
\[
\sin^2x=1-\cos^2x.
\]

Тогда
\[
2\cos^3x-2\cos x
=
1-\cos^2x,
\]
откуда
\[
2\cos^3x+\cos^2x-2\cos x-1=0.
\]

Делаем замену
\[
t=\cos x,
\qquad
-1\le t\le1.
\]

Получаем кубическое уравнение
\[
2t^3+t^2-2t-1=0.
\]

Группировка:
\[
(2t^3+t^2)+(-2t-1)=0,
\]
\[
t^2(2t+1)-(2t+1)=0,
\]
\[
(2t+1)(t^2-1)=0.
\]

Следовательно,
\[
t=-\frac12,
\qquad
t=-1,
\qquad
t=1.
\]

\begin{noteBox}
\textbf{Эвристика.}
Если в многочлене чётное число слагаемых, имеет смысл проверить возможность группировки.
Это полезный ориентир, который работает во многих задачах.
\end{noteBox}

Возвращаемся к косинусу:
\[
\cos x=-\frac12,
\qquad
\cos x=-1,
\qquad
\cos x=1.
\]

Для
\[
\cos x=-\frac12
\]
получаем две серии:
\[
x=\frac{2\pi}{3}+2\pi n,
\qquad
x=\frac{4\pi}{3}+2\pi n,
\qquad
n\in\Z.
\]

Ограничение
\[
\sin x\le0
\]
оставляет серию
\[
x=\frac{4\pi}{3}+2\pi n.
\]

Для
\[
\cos x=\pm1
\]
имеем
\[
\sin x=0,
\]
поэтому ограничение выполняется.

Две серии удобно объединить:
\[
x=\pi n,
\qquad
n\in\Z.
\]

Полный ответ:
\[
\boxed{
x=\pi n
\quad\text{или}\quad
x=\frac{4\pi}{3}+2\pi n,
\qquad
n\in\Z.
}
\]

\section{Общие формулы и табличные значения}

При
\[
|a|\le1
\]
для синуса:
\[
\sin x=a
\iff
\begin{cases}
x=\arcsin a+2\pi n,\\
x=\pi-\arcsin a+2\pi n,
\end{cases}
\qquad
n\in\Z.
\]

Для косинуса:
\[
\cos x=a
\iff
x=\pm\arccos a+2\pi n,
\qquad
n\in\Z.
\]

\begin{keybox}
Для значений
\[
0,\qquad 1,\qquad -1
\]
обычно быстрее пользоваться специальными случаями и единичной окружностью:
\[
\sin x=0
\Rightarrow
x=\pi n,
\]
\[
\cos x=0
\Rightarrow
x=\frac{\pi}{2}+\pi n,
\]
\[
\cos x=1
\Rightarrow
x=2\pi n,
\]
\[
\cos x=-1
\Rightarrow
x=\pi+2\pi n.
\]
\end{keybox}

\begin{table}[ht]
\centering
\renewcommand{\arraystretch}{1.28}

\caption{Базовые значения синуса и косинуса}

\begin{tabular}{@{}c*{5}{c}@{}}
\toprule

$x$
&
$0$
&
$\frac{\pi}{6}$
&
$\frac{\pi}{4}$
&
$\frac{\pi}{3}$
&
$\frac{\pi}{2}$
\\

\midrule

$\sin x$
&
$0$
&
$\frac12$
&
$\frac{\sqrt2}{2}$
&
$\frac{\sqrt3}{2}$
&
$1$
\\

$\cos x$
&
$1$
&
$\frac{\sqrt3}{2}$
&
$\frac{\sqrt2}{2}$
&
$\frac12$
&
$0$
\\

\bottomrule
\end{tabular}
\end{table}

Мнемоника для первой четверти:
\[
\sin:
\quad
\frac{\sqrt0}{2},
\quad
\frac{\sqrt1}{2},
\quad
\frac{\sqrt2}{2},
\quad
\frac{\sqrt3}{2},
\quad
\frac{\sqrt4}{2}.
\]

Для косинуса тот же ряд читается в обратном порядке.

\section{Отбор корней на промежутке}

Пусть нужно выбрать корни серии
\[
x=\frac{4\pi}{3}+2\pi n
\]
на отрезке
\[
\left[
-4\pi,
-\frac{5\pi}{2}
\right].
\]

Подставляем формулу корня в двойное неравенство:
\[
-4\pi
\le
\frac{4\pi}{3}+2\pi n
\le
-\frac{5\pi}{2}.
\]

Умножим все части на \(6\):
\[
-24\pi
\le
8\pi+12\pi n
\le
-15\pi.
\]

Вычтем \(8\pi\):
\[
-32\pi
\le
12\pi n
\le
-23\pi.
\]

Делим на положительное число \(12\pi\):
\[
-\frac83
\le
n
\le
-\frac{23}{12}.
\]

Единственное целое значение:
\[
n=-2.
\]

Подставляем его в исходную формулу серии:
\[
x
=
\frac{4\pi}{3}
+
2\pi(-2)
=
-\frac{8\pi}{3}.
\]

Для серии
\[
x=\pi n
\]
на том же отрезке:
\[
-4
\le
n
\le
-\frac52.
\]

Следовательно,
\[
n=-4,-3.
\]

Все корни разобранного уравнения на указанном отрезке:
\[
\boxed{
-4\pi,
\quad
-3\pi,
\quad
-\frac{8\pi}{3}.
}
\]

\begin{noteBox}
\textbf{Контроль при делении неравенства.}

При делении на отрицательное число направления знаков неравенства меняются.

В разобранном примере
\[
12\pi>0,
\]
поэтому направления знаков сохраняются.
\end{noteBox}

\section{Типичные ошибки}

\begin{itemize}

  \item
  Возвести
  \[
  \sqrt{F(x)}=G(x)
  \]
  в квадрат и потерять условие
  \[
  G(x)\ge0.
  \]

  \item
  Начать алгебраические преобразования, оставив разные аргументы тригонометрических функций.

  \item
  Подменить
  \[
  \sin^2x
  \]
  выражением
  \[
  1-\cos x
  \]
  вместо
  \[
  1-\cos^2x.
  \]

  \item
  При подстановке забыть скобки перед выражением со знаком минус.

  \item
  После замены
  \[
  t=\sin x
  \qquad\text{или}\qquad
  t=\cos x
  \]
  забыть ограничение
  \[
  -1\le t\le1.
  \]

  \item
  Для значения
  \[
  a\in(-1,1)
  \]
  записать только одну серию корней уравнения
  \[
  \sin x=a
  \]
  или
  \[
  \cos x=a.
  \]

  \item
  Потерять условие
  \[
  n\in\Z
  \]
  в общей формуле.

  \item
  При отборе корней подставить найденное целое \(n\) в промежуточное неравенство вместо исходной формулы серии.

  \item
  При делении двойного неравенства на отрицательное число забыть изменить направления обоих знаков.

\end{itemize}

% ============================================================
% Блок-схема
% ============================================================

\clearpage

\newgeometry{
  top=1.25cm,
  bottom=1.25cm,
  left=1.6cm,
  right=1.6cm
}

\thispagestyle{empty}

\begin{center}
{\LARGE
\bfseries
\color{darkaccent}
Алгоритм: тригонометрическое уравнение с радикалом}
\end{center}

\vspace{0.25cm}

\begin{center}

\begin{tikzpicture}[node distance=4.1mm]

\node[flowStart] (start)
{Уравнение с радикалом и тригонометрией};

\node[
  flowProcess,
  below=of start
] (radical)
{
Приведите к виду
$\sqrt{F(x)}=G(x)$:
запишите
$F(x)=G^2(x)$
и ограничение
$G(x)\ge0$
};

\node[
  flowProcess,
  below=of radical
] (angles)
{
Формулами приведения сведите аргументы к одному виду
};

\node[
  flowProcess,
  below=of angles
] (functions)
{
Через
$\sin^2x+\cos^2x=1$
сведите квадраты к одной функции
};

\node[
  flowProcess,
  below=of functions
] (subst)
{
Сделайте замену
$t=\sin x$
или
$t=\cos x$;
учтите
$-1\le t\le1$
};

\node[
  flowProcess,
  below=of subst
] (factor)
{
Решите алгебраическое уравнение и разложите выражение на множители, если это возможно
};

\node[
  flowProcess,
  below=of factor
] (restore)
{
Вернитесь к тригонометрии и выпишите все серии корней
};

\node[
  flowProcess,
  below=of restore
] (filter)
{
Проверьте серии по исходному ограничению
$G(x)\ge0$
};

\node[
  flowDecision,
  below=6mm of filter
] (interval)
{
Нужно отобрать корни на промежутке?
};

\node[
  flowProcess,
  below=7mm of interval
] (select)
{
Для каждой серии составьте двойное неравенство,
найдите целые $n$
и подставьте их в формулу корня
};

\node[
  flowStart,
  below=6mm of select
] (finish)
{
Полный ответ проверен
};

\draw[flowArrow]
(start) -- (radical);

\draw[flowArrow]
(radical) -- (angles);

\draw[flowArrow]
(angles) -- (functions);

\draw[flowArrow]
(functions) -- (subst);

\draw[flowArrow]
(subst) -- (factor);

\draw[flowArrow]
(factor) -- (restore);

\draw[flowArrow]
(restore) -- (filter);

\draw[flowArrow]
(filter) -- (interval);

\draw[flowArrow]
(interval)
--
node[right,font=\small]{да}
(select);

\draw[flowArrow]
(select) -- (finish);

\draw[flowArrow]
(interval.east)
--
++(2cm,0)
|-
node[
  pos=0.25,
  right,
  font=\small,
  align=left
]
{нет: оставить\\общие серии}
(finish.east);

\end{tikzpicture}

\end{center}

\clearpage
\restoregeometry

% ============================================================
% Тренировочный блок
% ============================================================

\section*{Тренировочный блок — Путь А}

\thispagestyle{fancy}




\subsection*{Задание 1}

Решите уравнение
\[
\sqrt{2\sin^3x-2\sin x}
=
\sin\left(
\frac{3\pi}{2}-x
\right).
\]

\begin{enumerate}[label=\alph*)]

  \item
  Найдите все решения уравнения.

  \item
  Отберите решения, принадлежащие отрезку
  \[
  [-3\pi,-\pi].
  \]

\end{enumerate}

\vspace{1em}

\textbf{Самопроверка перед сдачей:}

\begin{itemize}

  \item
  записано ограничение, появившееся при возведении в квадрат;

  \item
  использована формула приведения;

  \item
  после замены учтено
  \[
  -1\le t\le1;
  \]

  \item
  выписаны все серии тригонометрических корней;

  \item
  серии проверены по ограничению;

  \item
  при отборе найдено целое значение параметра каждой подходящей серии.

\end{itemize}

\vfill

\begin{center}
\small
\color{darkaccent}
Лёвин Артём Александрович эксклюзивно для Тимофея Васильченко
\end{center}

% ============================================================
% Ответы
% ============================================================

\clearpage

\section*{Ответы}

\renewcommand{\arraystretch}{1.45}

\begin{table}[ht]
\centering

\begin{tabularx}{\linewidth}{@{}p{2.1cm}X@{}}

\toprule

\textbf{Задание}
&
\textbf{Ответ}
\\

\midrule

1a
&
\(
\displaystyle
x=\frac{7\pi}{6}+2\pi n
\)
или
\(
\displaystyle
x=\frac{\pi}{2}+\pi n,
\)
где
\(
n\in\Z.
\)
\\

\addlinespace

1b
&
\(
\displaystyle
-\frac{17\pi}{6},
\quad
-\frac{5\pi}{2},
\quad
-\frac{3\pi}{2}.
\)
\\

\bottomrule

\end{tabularx}
\end{table}

\end{document}